\documentclass[11pt]{article}

\usepackage{../langcalc}
\usepackage{tikz}

\begin{document}

\courseheader
  {Day 02}
  {When Does a Sum Become an Area?}
  {LMTH 2040 --- Calculus: A Multidisciplinary Approach}
  {Monday, August 31, 2026}

\begin{keyidea}
Today we will begin with a pattern you already know and ask what happens
when we change its scale.

Our path will be

\[
\text{pattern}
\longrightarrow
\text{sum}
\longrightarrow
\text{picture}
\longrightarrow
\text{computation}
\longrightarrow
\text{area}.
\]

The central question is:

\begin{center}
\Large\bfseries
When does a sum become an area?
\end{center}
\end{keyidea}


% =========================================================
\section{1. Triangular Numbers Revisited}
% =========================================================

Recall the triangular numbers

\[
T_n = 1+2+3+\cdots+n.
\]

For example,

\[
T_1=1,
\qquad
T_2=3,
\qquad
T_3=6,
\qquad
T_4=10.
\]

\subsection*{Compare}

With your group, compare your solutions to the triangular-number
problem from Homework 01.

Identify at least two different ways of representing or determining
\(T_n\).

\vspace{0.15in}

\textbf{Representation 1:}

\vspace{0.65in}

\textbf{Representation 2:}

\vspace{0.65in}


\subsection*{What does each representation reveal?}

Complete the table.

\begin{center}
\begin{tabularx}{0.95\textwidth}{@{}p{0.24\textwidth}XX@{}}
\toprule
& \textbf{What is easy to see?}
& \textbf{What is difficult to see?}
\\
\midrule

Picture / diagram
&
\vspace{0.42in}
&
\\
\midrule

Sum
&
\vspace{0.42in}
&
\\
\midrule

Formula
&
\vspace{0.42in}
&
\\
\midrule

Computation
&
\vspace{0.42in}
&
\\

\bottomrule
\end{tabularx}
\end{center}


\begin{reflection}
A numerical pattern can provide \emph{evidence} for a conjecture.
A mathematical argument explains why the conjecture must continue
to hold.

Where does this distinction appear in your solution to the triangular
number problem?
\end{reflection}

\vspace{0.55in}


% =========================================================
\section{2. Changing the Scale}
% =========================================================

Suppose we place the \(n\)th triangular number inside an
\(n\times n\) square.

The square contains \(n^2\) unit squares.

The triangle contains \(T_n\) objects.

So consider the ratio

\[
A_n=\frac{T_n}{n^2}.
\]


\subsection*{Calculate}

Without using the general formula for \(T_n\), calculate:

\[
\begin{array}{c|c|c}
n & T_n & A_n=T_n/n^2\\
\hline
1 & &\\[0.12in]
2 & &\\[0.12in]
3 & &\\[0.12in]
4 & &\\[0.12in]
5 & &\\[0.12in]
10 & &
\end{array}
\]


\subsection*{Conjecture}

What appears to happen to \(A_n\) as \(n\) becomes large?

\vspace{0.55in}


\subsection*{Use the formula}

Now use

\[
T_n=\frac{n(n+1)}{2}
\]

to simplify

\[
A_n=\frac{T_n}{n^2}.
\]

\vspace{0.65in}

What does your expression suggest happens as \(n\) becomes very large?

\vspace{0.5in}


\begin{reflection}
We have changed the question.

Instead of asking how large \(T_n\) becomes, we are asking how large it
is \emph{relative to the square containing it}.

Why might changing the scale reveal a pattern that is difficult to see
in the original numbers?
\end{reflection}


% =========================================================
\newpage
\section{3. Turn the Sum into a Picture}
% =========================================================

Start again with

\[
A_n
=
\frac{T_n}{n^2}
=
\frac{1}{n^2}
\left(
1+2+\cdots+n
\right).
\]

Rewrite this expression by distributing the two factors of \(n\):

\[
A_n
=
\frac1n
\left(
\frac1n
+
\frac2n
+
\frac3n
+
\cdots
+
\frac nn
\right).
\]


\subsection*{A geometric interpretation}

Consider a square with side length \(1\).

Divide the horizontal axis into \(n\) equal pieces.

Each piece therefore has width

\[
\boxed{\phantom{\frac{1}{n}}}.
\]

Now interpret

\[
\frac1n,\quad
\frac2n,\quad
\frac3n,\quad
\ldots,\quad
\frac nn
\]

as heights.


\begin{center}
\begin{tikzpicture}[scale=6]

    % axes
    \draw[->] (0,0) -- (1.08,0) node[right] {$x$};
    \draw[->] (0,0) -- (0,1.08) node[above] {$y$};

    % unit square guides
    \draw[dashed] (1,0) -- (1,1);
    \draw[dashed] (0,1) -- (1,1);

    % labels
    \node[below] at (1,0) {$1$};
    \node[left] at (0,1) {$1$};

\end{tikzpicture}
\end{center}


For \(n=4\), draw four rectangles with width \(1/4\) and heights

\[
\frac14,\qquad
\frac24,\qquad
\frac34,\qquad
\frac44.
\]


\subsection*{Area}

What is the area of the first rectangle?

\[
\text{area}
=
\boxed{\phantom{\frac14}}
\cdot
\boxed{\phantom{\frac14}}
=
\boxed{\phantom{\frac1{16}}}
\]

What is the area of the second?

\[
\text{area}
=
\boxed{\phantom{\frac14}}
\cdot
\boxed{\phantom{\frac24}}
=
\boxed{\phantom{\frac2{16}}}
\]

What is the total area of all four rectangles?

\vspace{0.45in}

How is this related to \(T_4/4^2\)?

\vspace{0.45in}


% =========================================================
\section{4. Refine the Picture}
% =========================================================

Draw the same construction for \(n=8\).

\begin{center}
\begin{tikzpicture}[scale=6]

    \draw[->] (0,0) -- (1.08,0) node[right] {$x$};
    \draw[->] (0,0) -- (0,1.08) node[above] {$y$};

    \draw[dashed] (1,0) -- (1,1);
    \draw[dashed] (0,1) -- (1,1);

    \foreach \x in {0.125,0.25,0.375,0.5,0.625,0.75,0.875}
        \draw[very thin] (\x,0) -- (\x,0.025);

    \node[below] at (1,0) {$1$};
    \node[left] at (0,1) {$1$};

\end{tikzpicture}
\end{center}


What changes when we move from \(n=4\) to \(n=8\)?

\vspace{0.45in}

What stays the same?

\vspace{0.45in}

What do you expect the picture to look like for \(n=100\)?

\vspace{0.45in}


\begin{problem}
Imagine continuing this process:

\[
n=4,\qquad
n=8,\qquad
n=100,\qquad
n=1000,\qquad
\ldots
\]

What geometric object appears to emerge from the rectangles?

Draw it on one of the diagrams above.
\end{problem}


% =========================================================
\newpage
\section{5. Let the Spreadsheet Do the Repetition}
% =========================================================

Open the starter spreadsheet.

The first column contains values of \(n\). Your task is to make the
spreadsheet investigate the triangular numbers.


\subsection*{Method A: Build the sequence}

Triangular numbers satisfy

\[
T_n=T_{n-1}+n.
\]

Use this relationship to construct the sequence in the spreadsheet.

What formula did you enter?

\vspace{0.35in}

What information does the spreadsheet need in order to begin this
process?

\vspace{0.35in}


\subsection*{Method B: Calculate directly}

We also know

\[
T_n=\frac{n(n+1)}2.
\]

Use this formula in a new column.

Do the two methods produce the same values?

\[
\boxed{\text{Yes}}
\qquad
\boxed{\text{No}}
\]

If they agree, does that prove that the formula is correct for every
positive integer \(n\)? Explain.

\vspace{0.6in}


\subsection*{Normalize}

Create another column containing

\[
A_n=\frac{T_n}{n^2}.
\]

Investigate at least

\[
n=10,\qquad
100,\qquad
1000.
\]

Record your results.

\[
\begin{array}{c|c}
n&A_n\\
\hline
10&\\[0.15in]
100&\\[0.15in]
1000&
\end{array}
\]


What value does the spreadsheet suggest \(A_n\) is approaching?

\vspace{0.35in}

How confident are you?

\vspace{0.35in}

What evidence would make you more confident?

\vspace{0.4in}


\begin{reflection}
The spreadsheet can calculate thousands of examples very quickly.

What can the spreadsheet tell us?

\vspace{0.25in}

What can it \emph{not} tell us by calculation alone?
\end{reflection}


% =========================================================
\section{6. A New Notation}
% =========================================================

Instead of writing

\[
1+2+3+\cdots+n,
\]

mathematicians often write

\[
\sum_{k=1}^{n} k.
\]

The symbol

\[
\sum
\]

means that we are adding a collection of terms.

So

\[
T_n
=
\sum_{k=1}^{n}k.
\]


Rewrite each expression using sigma notation.

\[
1+2+3+\cdots+100
\]

\vspace{0.25in}

\[
1^2+2^2+3^2+\cdots+n^2
\]

\vspace{0.25in}

\[
\frac1n+\frac2n+\frac3n+\cdots+\frac nn
\]

\vspace{0.35in}


Now return to our rectangle construction.

Each rectangle has width \(1/n\), while its height is \(k/n\).

Its area is therefore

\[
\frac{k}{n}\frac1n.
\]

The total area of all the rectangles is

\[
\boxed{
\sum_{k=1}^{n}
\frac{k}{n}\frac1n
}.
\]


% =========================================================
\newpage
\section{7. When Does a Sum Become an Area?}
% =========================================================

The top edges of our rectangles approach the line

\[
y=x.
\]

Draw the line \(y=x\) and several rectangles below.

\begin{center}
\begin{tikzpicture}[scale=7]

    \draw[->] (0,0) -- (1.08,0) node[right] {$x$};
    \draw[->] (0,0) -- (0,1.08) node[above] {$y$};

    \draw[dashed] (1,0) -- (1,1);
    \draw[dashed] (0,1) -- (1,1);

    \node[below] at (1,0) {$1$};
    \node[left] at (0,1) {$1$};

\end{tikzpicture}
\end{center}


As \(n\) becomes larger:

\begin{itemize}[leftmargin=*]

\item the rectangles become \rule{1.2in}{0.4pt};

\item the number of rectangles becomes \rule{1.2in}{0.4pt};

\item their total area appears to approach \rule{1.2in}{0.4pt}.

\end{itemize}


We can describe the process symbolically as

\[
\lim_{n\to\infty}
\sum_{k=1}^{n}
\frac{k}{n}\frac1n
=
\frac12.
\]


\begin{keyidea}
The expression on the left is a sequence of \emph{finite sums}.

Every individual calculation uses only finitely many rectangles.

Yet the limiting process appears to determine an exact area:

\[
\boxed{\frac12}.
\]

How can an infinite process produce an exact finite quantity?
\end{keyidea}


\subsection*{One more piece of notation}

Modern calculus gives the limiting area another name:

\[
\boxed{
\int_0^1 x\,dx=\frac12
}
\]

For now, do not worry about calculating integrals using rules or
formulas.

Instead, read the symbol as a question:

\begin{center}
\large
\emph{What quantity is produced when increasingly fine pieces
are accumulated?}
\end{center}


% =========================================================
\section{8. Looking Back / Looking Forward}
% =========================================================

Return to today's path:

\[
\text{pattern}
\longrightarrow
\text{sum}
\longrightarrow
\text{picture}
\longrightarrow
\text{computation}
\longrightarrow
\text{area}.
\]

Choose one transition in this chain.

What new idea became visible when we changed representations?

\vspace{0.65in}


\begin{historicalnote}
Today we used algebraic formulas, decimal calculations, sigma notation,
a spreadsheet, and the modern language of limits to investigate an area.

But the problem is much older than any of these tools.

On Wednesday we will ask how Archimedes could establish exact results
about curved figures using only geometry, finite constructions, and
carefully controlled approximation.
\end{historicalnote}


\subsection*{Exit Question}

Suppose someone says:

\begin{quote}
``A computer can calculate this sum for an enormous value of \(n\),
so we do not need the limit.''
\end{quote}

What would you say in response?

\vspace{0.8in}

\clearpage



% =========================================================
\section*{For Wednesday, September 2}
% =========================================================

\courselabel{Reading}

For our next class, read the selections from:

\begin{itemize}[leftmargin=*]

    \item \textbf{Archimedes}, \emph{The Quadrature of the Parabola}.

    \item \textbf{Michel Serres}, \emph{The Birth of Physics},
    ``The Work of Archimedes,'' together with the short selection
    ``Archimedes on Thinking Deviation.''

\end{itemize}

As you read, do not worry about understanding every step of
Archimedes' argument or every claim made by Serres. Come to class
prepared to discuss the following questions:

\begin{enumerate}[leftmargin=*]

    \item What problem is Archimedes trying to solve?

    \item What role do approximation, geometry, and infinity appear
    to play in his argument?

    \item What does Serres see in Archimedes that might not be
    apparent from reading the mathematics alone?

\end{enumerate}

\begin{reflection}
Choose one passage, diagram, or mathematical step from either reading
that you would like us to discuss in class. Mark it and bring it with
you on Wednesday.

No written response is due.
\end{reflection}

\end{document}